
This section turns the model's maintained link into a portfolio-risk statement.
It first asks what exposure separation each observed pair can certify, then
aggregates those pairwise floors under the coherent joint law.
The resulting credit lowers the benchmark in which no cross-asset separation
can be certified.

\subsection{Observable pairwise floors}

The carrier and slack restrictions first yield the minimum exposure separation
supported by each pair's observed information distance.
For assets $i$ and $j$, define
\begin{equation}\label{eq:p3-information-floor}
  \ell_{ij}
  :=\left[
    L^{-1}W_2(C_i,C_j)-\tau_i-\tau_j
  \right]_+,
  \qquad [a]_+:=\max\{a,0\}.
\end{equation}
The antilipschitz carrier first retains at least the fraction $L^{-1}$ of the
observed W2 separation.
The triangle inequality then deducts the two slack radii.
Truncation at zero preserves the nonnegative lower bound, giving
\begin{equation}\label{eq:p3-floor-validity}
  \ell_{ij}\le W_2(P_i,P_j).
\end{equation}
Every term in \eqref{eq:p3-information-floor} therefore has a distinct role:
$W_2(C_i,C_j)$ is observed geometry, $L$ controls common distortion, and the
$\tau_i$ terms absorb asset-specific departures from the common carrier.
In the empirical application, $W_2(C_i,C_j)$ is the least root-mean-square
displacement required to align the two firms' information distributions.
A positive $\ell_{ij}$ rules out exposure laws that are closer than this floor;
a zero floor means only that the maintained restrictions certify no positive
separation for that pair.

The candidate portfolio credit weights each pairwise floor by the extent to
which both assets enter a normalized long-only portfolio:
\begin{equation}\label{eq:p3-certificate-credit}
  \mathcal C(q)
  :=\frac12\sum_i \sum_j q_i q_j\ell_{ij}^2.
\end{equation}
The factor one half compensates for counting both ordered pairs $(i,j)$ and
$(j,i)$.
Because diagonal distances are zero, this is also
$\sum_{i<j} q_i q_j\ell_{ij}^2$.
The ordered-pair sum is portfolio-specific: separation between two assets
contributes only to the extent that both are held.
Conditional on the declared carrier and slack parameters, $\mathcal C(q)$ is
computed from observed information geometry rather than a return covariance
matrix.

\subsection{From floors to coherent portfolio variance}

Pairwise floors constrain individual asset pairs, but portfolio risk is a joint
object.
To aggregate the floors without combining incompatible pairwise couplings,
apply the weighted Hilbert-space polarization identity under the same joint law
$J$:
\begin{equation}\label{eq:p3-weighted-polarization}
  \left\lVert\sum_i q_i B_i\right\rVert^2
  =\sum_i q_i\lVert B_i\rVert^2
  -\frac12\sum_i \sum_j q_i q_j\lVert B_i-B_j\rVert^2.
\end{equation}
Taking expectations under the coherent law $J$ is valid when the pairwise
inner products are integrable.
For each pair, the realized $J$-marginal is an admissible coupling of
$P_i$ and $P_j$.
Its expected squared displacement is therefore no smaller than
$W_2^2(P_i,P_j)$, which in turn is no smaller than $\ell_{ij}^2$ by
\eqref{eq:p3-floor-validity}.
Substituting those lower bounds into the subtracted term of
\eqref{eq:p3-weighted-polarization} yields the main result.

\begin{theorem}[Information-certified coherent portfolio variance]
  \label{thm:p3-information-certified-cap}
  Suppose the common carrier satisfies \eqref{eq:p3-antilipschitz}, the
  asset-specific maps satisfy \eqref{eq:p3-synchronous-slack}, the exposure
  laws are marginals of one joint law $J$, and all pairwise inner products are
  integrable.
  Then every normalized long-only portfolio obeys
  \begin{equation}\label{eq:p3-systematic-cap}
    V_{\mathrm{sys}}(q)
    \le \sum_i q_i v_i-\mathcal C(q).
  \end{equation}
\end{theorem}

The theorem is an upper bound on systematic portfolio variance, not a point
estimate of a covariance matrix.
Its benchmark $\sum_i q_i v_i$ is the weighted marginal second-moment term that
would remain if no cross-asset separation could be certified.
The observable characteristic laws deduct $\mathcal C(q)$ from that benchmark.
Larger distances tighten the deduction; larger distortion or slack weakens it.
In finance terms, $\mathcal C(q)$ is the diversification relief that the
observed information can certify under the maintained transmission model.

The result is conservative in two specific ways, which the next subsection
makes precise and addresses jointly.
Pairwise optimal couplings need not all coexist inside one joint law, so the
weighted pairwise floor understates the separation any coherent joint exposure
law must incur.
Equation \eqref{eq:p3-certificate-credit} also deducts each pair's two slack
radii separately and truncates the result pair by pair, discarding every pair
whose observed separation falls below its combined radius.

\subsection{The sharp multi-firm certificate}\label{sec:sharp-certificate}

The pairwise credit relaxes a single multi-firm object.
That object measures the least total separation compatible with all observed
characteristic laws under one common coupling.
We refer to it as weighted multi-firm transport dispersion; for observed
characteristic laws, the formal object below is weighted characteristic
dispersion.
Stating it directly tightens the deduction, and its two-asset case is exactly
the pairwise covariance envelope of \citet{gawronsky_continuous_2026}, so the
sharpening does not introduce a second theory.

Fix normalized risk weights $q$ as above.
Within $\mathcal D_q$, these weights are inputs: the infimum varies the coherent
common coupling while holding $q$ fixed.
Only the portfolio-choice problem in \Cref{sec:portfolio-choice} later treats
$q$ as a decision variable.

\begin{definition}[Weighted characteristic dispersion]
  \label{def:p3-dispersion}
  Let $\Pi$ denote the set of joint laws on $\mathcal X^n$ whose $i$th marginal
  is $C_i$. For characteristic laws $C_1,\dots,C_n$ on
  $(\mathcal X,d_{\mathcal X})$ and $(X_1,\dots,X_n)\sim\gamma$,
  \begin{equation}\label{eq:p3-dispersion}
    \mathcal D_q(C_1,\dots,C_n)
    :=\inf_{\gamma\in\Pi}
    \E_\gamma\left[\sum_{i<j}q_i\,q_j\,d_{\mathcal X}(X_i,X_j)^2\right].
  \end{equation}
\end{definition}

The infimum is what makes the quantity conservative. It asks how close the
laws could be under their most favourable common coupling, not how separated
one selected matching makes them appear. Its units are squared characteristic
distance, and its value depends on the ground metric and on $q$.

Two facts locate \eqref{eq:p3-dispersion} against objects already in play.
First, restricting a common coupling to any pair leaves an admissible coupling
of that pair, whose expected squared displacement is therefore at least
$W_2^2(C_i,C_j)$. Summing gives
\begin{equation}\label{eq:p3-pairwise-relaxation}
  \mathcal D_q(C_1,\dots,C_n)
  \ \ge\ \sum_{i<j}q_i\,q_j\,W_2^2(C_i,C_j).
\end{equation}
The weighted pairwise sum is thus a computable lower bound on the multi-firm
object, which is the inequality \eqref{eq:p3-certificate-credit} exploits.
It can be built once from pairwise distances, but it does not retain the full
restriction that all asset marginals share one coupling.
Second, on a Hilbert space the same value has a free-centre Wasserstein
barycentre representation: for exposure laws with finite second moments,
\begin{equation}\label{eq:p3-barycenter-representation}
  \mathcal D_q(P_1,\dots,P_n)
  =\inf_{Q}\ \sum_i q_i\,W_2^2(P_i,Q),
\end{equation}
where $Q$ ranges over finite-second-moment laws on $\mathcal H$.
For fixed $q$, only $Q$ varies.
This shares the barycentric geometry of \citet{gawronsky_spatial_2027} but is not
its target-anchored Wasserstein barycentric reconstruction; that problem fixes
$i$ and its alignments and varies $w_i$ to form $W^\flat$.
\Cref{prop:p3-barycenter-representation} states this in the appendix. Only
equality of the two infimum values is used; neither an optimal barycentre nor
an optimal joint plan is assumed to exist.

The two-asset case recovers the pairwise foundation exactly.

\begin{proposition}[Two-asset nesting]\label{thm:p3-pairwise-bridge}
  Let $n=2$ and suppose some coherent joint law attains the dispersion minimum
  for $P_1,P_2$. Then the systematic portfolio variance it attains is
  \begin{equation}\label{eq:p3-two-asset-nesting}
    q_1v_1+q_2v_2-q_1q_2\,W_2^2(P_1,P_2).
  \end{equation}
\end{proposition}
Substituting $n=2$ into the dispersion definition gives
$\mathcal D_q(P_1,P_2)=q_1q_2W_2^2(P_1,P_2)$.
Weighted polarization then yields \eqref{eq:p3-two-asset-nesting}.

Equation \eqref{eq:p3-two-asset-nesting} is the covariance-envelope correction
of \citet{gawronsky_continuous_2026} written in portfolio form.
The multi-firm construction extends that pairwise foundation to a coherent
joint exposure law.

The dispersion in \eqref{eq:p3-dispersion} remains measured in observed
characteristic units, whereas the variance bound concerns latent exposure
units.
The same carrier and slack restrictions transfer the whole multi-firm object
between these spaces rather than treating each pair separately.
Aggregate the firm-specific mismatch through the weighted root-mean-square
slack radius
\begin{equation}\label{eq:p3-rms-slack}
  \tau_q:=\Big(\sum_i q_i\tau_i^2\Big)^{1/2}.
\end{equation}

\begin{theorem}[Characteristic-to-exposure dispersion]
  \label{thm:p3-aggregate-transfer}
  Suppose the common carrier satisfies \eqref{eq:p3-antilipschitz} and the
  asset-specific maps satisfy \eqref{eq:p3-synchronous-slack}. Then
  \begin{equation}\label{eq:p3-aggregate-transfer}
    \sqrt{\mathcal D_q(P_1,\dots,P_n)}
    \ \ge\
    \left[
      L^{-1}\sqrt{\mathcal D_q(C_1,\dots,C_n)}-\tau_q
    \right]_+ .
  \end{equation}
\end{theorem}

For the transfer proof, write
$T_i:=t_\#C_i$ and abbreviate
$\mathcal D_q(C_1,\ldots,C_n)$,
$\mathcal D_q(P_1,\ldots,P_n)$, and
$\mathcal D_q(T_1,\ldots,T_n)$ by $\mathcal D_q(C)$,
$\mathcal D_q(P)$, and $\mathcal D_q(T)$.
The antilipschitz inequality implies, after pushing any coupling of the
$C_i$ through $t$,
\[
  \sqrt{\mathcal D_q(C)}
  \le L\sqrt{\mathcal D_q(T)}.
\]
The synchronous coupling $(u_i(X_i),t(X_i))$ gives
\[
  W_2(P_i,T_i)\le\tau_i.
\]
The weighted product-space stability inequality then yields
\[
  \left|\sqrt{\mathcal D_q(P)}-\sqrt{\mathcal D_q(T)}\right|
  \le\left(\sum_i q_i W_2^2(P_i,T_i)\right)^{1/2}
  \le\tau_q.
\]
Combining the last two displays gives
\[
  \sqrt{\mathcal D_q(P)}
  \ge L^{-1}\sqrt{\mathcal D_q(C)}-\tau_q.
\]
Both sides of the left-hand inequality are nonnegative, so taking the
positive part proves \eqref{eq:p3-aggregate-transfer}.
Thus, even under the most favourable common coupling of the observed laws, the
carrier preserves a minimum amount of latent exposure dispersion after one
aggregate allowance for firm-specific mismatch.

Squaring this transferred floor gives the sharp observable credit
\begin{equation}\label{eq:p3-sharp-credit}
  \mathcal C^{\sharp}(q)
  :=\left[
    L^{-1}\sqrt{\mathcal D_q(C_1,\dots,C_n)}-\tau_q
  \right]_+^{2}.
\end{equation}

\begin{theorem}[Sharp information-certified portfolio variance]
  \label{thm:p3-sharp-certificate}
  Under the hypotheses of \Cref{thm:p3-information-certified-cap}, every
  normalized long-only portfolio obeys
  \begin{equation}\label{eq:p3-sharp-cap}
    V_{\mathrm{sys}}(q)\le\sum_i q_i\,v_i-\mathcal C^{\sharp}(q).
  \end{equation}
\end{theorem}

The proof is the following chain, where the first inequality uses the
dispersion infimum and the second uses
\Cref{thm:p3-aggregate-transfer}:
\[
  \begin{aligned}
    V_{\mathrm{sys}}(q)
    &=\sum_i q_i v_i-
    \E_J\!\left[\sum_{i<j} q_i q_j\lVert B_i-B_j\rVert^2\right]\\
    &\le\sum_i q_i v_i-\mathcal D_q(P)\\
    &\le\sum_i q_i v_i-\mathcal C^{\sharp}(q).
  \end{aligned}
\]
No attaining plan is needed for this bound.
The theorem therefore tightens the systematic-risk cap by preserving joint
compatibility across all assets instead of replacing the infimum over common
couplings with a sum of pairwise infima.

\begin{proposition}[Envelope exactness]\label{thm:p3-envelope-exactness}
  If some coherent joint law attains the dispersion minimum for
  $P_1,\dots,P_n$, then $\sum_i q_i\,v_i-\mathcal D_q(P_1,\dots,P_n)$ is the
  greatest systematic portfolio variance attainable across coherent joint laws.
\end{proposition}
An attaining joint law realizes $\mathcal D_q(P)$, while every other coherent
law has dispersion at least that infimum.
Consequently, none can have greater systematic variance than the value stated
in the proposition.
This exactness concerns the variance envelope generated by the fixed exposure
marginals and an attaining coherent law; it does not identify a return
covariance matrix from text.

\Cref{thm:p3-sharp-certificate} dominates
\Cref{thm:p3-information-certified-cap} in two independent respects.
Equation \eqref{eq:p3-pairwise-relaxation} makes the multi-firm dispersion at
least the weighted pairwise sum, and \eqref{eq:p3-sharp-credit} subtracts one
aggregate radius where \eqref{eq:p3-information-floor} subtracts
$\tau_i+\tau_j$ from every pair. Under a common radius $\tau$ the second effect
alone replaces a deduction of $2\tau$ per pair by a single $\tau$. Neither
improvement needs new data: both use the same observed W2 geometry and the same
declared $(L,\tau)$.

The direction of the improvement is substantive.
A larger credit is a tighter variance bound, and the sharp form replaces the
pairwise relaxation with the coherent multi-firm dispersion while applying
slack once at the aggregate level.
The pairwise credit remains the decision form because it is a quadratic
function of $q$ built once from the observed distance matrix, whereas
evaluating $\mathcal D_q$ at a candidate $q$ requires solving an inner
free-centre barycentre problem.
The later portfolio search varies $q$ outside those fixed-$q$ evaluations; it
is not part of the definition of $\mathcal D_q$.
The empirical exercise therefore reports the pairwise relaxation while the
sharp theorem records the coherent multi-firm benchmark against which that
computational relaxation is judged.

\subsection{Return bridge and residual sensitivity}

The results so far bound latent systematic exposure risk, not observed return
variance.
To obtain a total-return statement, standardize each asset to unit marginal
variance and decompose its return into systematic and residual components.
This return bridge is a maintained restriction rather than a consequence of
the information geometry.

\begin{assumption}[Return bridge]\label{ass:p3-return-bridge}
  Each standardized return splits as $R_i=S_i+\varepsilon_i$ into a scalar
  systematic component $S_i$ carrying the exposure geometry and a residual
  $\varepsilon_i$.
  Write $v_i=\operatorname{Var}(S_i)$.
  The residual vector is cross-orthogonal to the systematic component:
  \[
    \operatorname{Cov}(S_i,\varepsilon_j)=0
    \quad\text{for every }i\text{ and }j.
  \]
  For risk weights $q$, set $S_q:=\sum_i q_i S_i$ and write
  \[
    \operatorname{Var}(R_q)=\operatorname{Var}(S_q)+\mathcal R(q),
    \qquad
    \mathcal R(q):=\sum_i\sum_j
    q_i q_j\operatorname{Cov}(\varepsilon_i,\varepsilon_j).
  \]
  Unit marginal return variance and zero systematic--residual covariance give
  $\operatorname{Var}(\varepsilon_i)=1-v_i$.
  Zero cross-residual covariance is the benchmark; more generally, the excess
  residual contribution is bounded by a declared $\delta\ge0$.
\end{assumption}

Under zero cross-residual covariance, the systematic certificate and the
diagonal residual terms give
\begin{equation}\label{eq:p3-standardized-cap}
  \begin{aligned}
    \operatorname{Var}(R_q)
    &\le \sum_i q_i v_i-\mathcal C(q)
    +\sum_i q_i^2(1-v_i)\\
    &=1-\mathcal C(q)-\sum_i q_i(1-q_i)(1-v_i)\\
    &\le 1-\mathcal C(q).
  \end{aligned}
\end{equation}
The first line combines the systematic certificate with diagonal residual
variance.
The equality isolates the additional diversification relief supplied by
asset-specific residual risk.
Long-only normalization gives $q_i(1-q_i)\ge0$, so the final inequality
discards that nonnegative relief and leaves the simpler
$1-\mathcal C(q)$ certificate.
With residual contamination, retain the sensitivity wrapper
\begin{equation}\label{eq:p3-residual-budget-cap}
  \operatorname{Var}(R_q)\le1-\mathcal C(q)+\delta.
\end{equation}
Here $\delta$ is a portfolio residual-covariance budget used for sensitivity
analysis, not an estimated structural parameter.
To express the standardized result as a capital-allocation bound, separate
capital weights from shares of the perfect-dependence volatility budget.
For long-only capital weights $x_i\ge0$ with $\sum_i x_i=1$ and positive
marginal scales $\sigma_i$, define
\begin{equation}\label{eq:p3-risk-weights}
  A(x):=\sum_i x_i\sigma_i,
  \qquad
  q_i(x):=\frac{x_i\sigma_i}{A(x)}.
\end{equation}
Here $A(x)$ is the portfolio volatility under perfect positive dependence.
The normalized weight $q_i(x)$ is asset $i$'s share of that benchmark.
Thus $x_i$ are capital weights whereas $q_i$ are risk-budget shares; in
particular, $x_i\propto1/\sigma_i$ implies $q_i=1/n$.
Multiplying \eqref{eq:p3-residual-budget-cap} by $A(x)^2$ yields
\begin{equation}\label{eq:p3-raw-cap}
  \operatorname{Var}(R_x)
  \le
  A(x)^2\left\{1-\mathcal C(q(x))+\delta\right\}.
\end{equation}
This is the financial interpretation of the certificate:
$A(x)^2$ is the perfect-dependence benchmark, while
$A(x)^2\mathcal C(q(x))$ is the amount of raw systematic variance
ruled out by the
observed information geometry under the maintained transmission model.
The additive term $A(x)^2\delta$ makes residual sensitivity explicit.
Equation \eqref{eq:p3-raw-cap} is therefore a decision input rather than an
estimated covariance model.
At the benchmark $\delta=0$, \Cref{sec:portfolio-choice} turns this certified
risk bound into a portfolio rule by minimizing it over a feasible allocation
set.
